\section*{Chapitre \ref{ch:b1:poly} --- Polynômes}

\begin{solution}{exo:b1:poly:1}
$X^5 - 1 = (X^2 + X + 1)(X^3 - X^2 + 1) + (-X - 2)$. Étapes~: on
retranche $X^3 B$, puis $-X^2 B$, puis $B$~; le reste $-X - 2$ est de
degré $1 < 2$. \emph{Vérification en $X = 1$~:} $\;0 = 3 \times 1 +
(-3)$.

$2X^4 + X^3 - X + 3 = (X^2 - 2)(2X^2 + X + 4) + (X + 11)$.
\emph{Vérification en $X = 0$~:} $\;3 = (-2)(4) + 11$.
\end{solution}

\begin{solution}{exo:b1:poly:2}
$X^2 + X + 1 = (X - j)(X - j^2)$ avec $j = \eu^{2\iu\pi/3}$, $j^3 =
1$. Ce \omterm{def:b1:poly:def}{polynôme} \omterm{def:b1:arith:divides}{divise} $Q_n = X^{2n} + X^n + 1$ si et seulement si $j$
et $j^2$ sont racines de $Q_n$~; comme $Q_n$ est à coefficients réels,
$Q_n(j^2) = \conj{Q_n(j)}$, si bien que la condition se réduit à
$Q_n(j) = 0$. Or $Q_n(j) = j^{2n} + j^n + 1$ ne dépend que de $n$
modulo $3$~:
\begin{itemize}
  \item $n \equiv 0$~: $Q_n(j) = 1 + 1 + 1 = 3 \neq 0$~;
  \item $n \equiv 1$~: $Q_n(j) = j^2 + j + 1 = 0$~;
  \item $n \equiv 2$~: $Q_n(j) = j^4 + j^2 + 1 = j + j^2 + 1 = 0$.
\end{itemize}
Ainsi $X^2 + X + 1 \mid X^{2n} + X^n + 1$ exactement lorsque $3 \nmid
n$.
\end{solution}

\begin{solution}{exo:b1:poly:3}
D'après la \cref{prop:b1:poly:multiplicity}, $(X-1)^2 \mid P$ si et
seulement si $P(1) = P'(1) = 0$~:
\[
P(1) = 2 + a + b = 0, \qquad P'(1) = 4 + 3a + 2b = 0 .
\]
En résolvant~: $b = -a - 2$ et $4 + 3a - 2a - 4 = a = 0$, donc $a = 0$,
$b = -2$~: $P = X^4 - 2X^2 + 1 = (X^2 - 1)^2 = (X-1)^2 (X+1)^2$, ce qui
est la factorisation réelle.
\end{solution}

\begin{solution}{exo:b1:poly:4}
$X^3 - 1 = (X - 1)(X - j)(X - j^2)$ sur $\C$ ($j = \eu^{2\iu\pi/3}$),
et $(X - 1)(X^2 + X + 1)$ sur $\R$.

$X^4 + X^2 + 1 = (X^2 + X + 1)(X^2 - X + 1)$ sur $\R$ (développer, ou
remarquer que $X^4 + X^2 + 1 = (X^2+1)^2 - X^2$)~; sur $\C$, chaque
trinôme se scinde~: racines $j, j^2$ et $-j, -j^2$, c'est-à-dire
$\eu^{\pm 2\iu\pi/3}, \eu^{\pm\iu\pi/3}$.

$X^6 - 1 = \prod_{k=0}^{5} (X - \eu^{\iu k\pi/3})$ sur $\C$, et sur
$\R$~:
\[
X^6 - 1 = (X-1)(X+1)(X^2 + X + 1)(X^2 - X + 1),
\]
en groupant les paires \omterm{def:b1:complex:field}{conjuguées} $\eu^{\pm 2\iu\pi/3}$ et $\eu^{\pm
\iu\pi/3}$.
\end{solution}

\begin{solution}{exo:b1:poly:5}
\begin{enumerate}
  \item Soit $p/q$ (sous forme irréductible) une racine du \omterm{def:b1:poly:def}{polynôme
        unitaire} à coefficients entiers $X^3 + \dots + c_0$~: en
        chassant les dénominateurs dans $P(p/q) = 0$, on obtient $p^3
        = -q\,(\text{entier})$, donc $q \mid p^3$~; la primalité entre
        eux impose $q = \pm 1$~: la racine est un entier $p$, et $p
        \mid c_0$ (isoler $c_0$). Ici les candidats \omterm{def:b1:arith:divides}{divisent} $6$~: en
        testant, $P(1) = 0$, $P(2) = 0$, $P(3) = 0$. Donc $P =
        (X-1)(X-2)(X-3)$.
  \item Viète~: $s_1 = 6$, $s_2 = 11$, $s_3 = 6$. Somme des carrés~:
        $s_1^2 - 2s_2 = 36 - 22 = 14 = 1 + 4 + 9$, comme prévu. Somme
        des inverses~: $\frac{s_2}{s_3} = \frac{11}{6} = 1 + \frac12 +
        \frac13$, comme prévu.
\end{enumerate}
\end{solution}

\begin{solution}{exo:b1:poly:6}
Première étape de division de l'\omterm{met:b1:arith:euclid}{algorithme d'Euclide}~:
\[
(X + 1)(X^3 - X^2 + X - 1)
= X^4 - X^3 + X^2 - X + X^3 - X^2 + X - 1 = X^4 - 1 ,
\]
donc la division de $X^4 - 1$ par $X^3 - X^2 + X - 1$ est exacte
(quotient $X + 1$, reste $0$), et l'algorithme s'arrête aussitôt~:
\[
\gcd(X^4 - 1,\; X^3 - X^2 + X - 1) = X^3 - X^2 + X - 1
\]
(déjà \omterm{def:b1:poly:def}{unitaire}). La relation de Bézout est la relation triviale~:
$\gcd = 0 \cdot (X^4 - 1) + 1 \cdot (X^3 - X^2 + X - 1)$. Vérification
de cohérence par factorisation~: $X^3 - X^2 + X - 1 = (X - 1)(X^2 +
1)$, qui est bien le produit des facteurs irréductibles communs à $X^4
- 1 = (X-1)(X+1)(X^2+1)$.
\end{solution}

\begin{solution}{exo:b1:poly:7}
Comme $P \geq 0$ sur $\R$, ses racines réelles sont de \omterm{def:b1:poly:derivative}{multiplicité}
paire (en une racine de \omterm{def:b1:poly:derivative}{multiplicité} impaire, $P$ change de signe). En
utilisant \cref{cor:b1:poly:factorization} et en appariant, écrivons
\[
P = c \prod_i (X - a_i)^{2k_i} \prod_j \bigl((X - z_j)(X - \conj
z_j)\bigr)^{n_j},
\]
avec $c > 0$ (comportement en $+\infty$). Posons
\[
S = \sqrt c\, \prod_i (X - a_i)^{k_i} \prod_j (X - z_j)^{n_j}
\in \C[X],
\]
de sorte que $P = S\,\conj S$, où $\conj S$ a les coefficients
\omterm{def:b1:complex:field}{conjugués}. Décomposons $S = A + \iu B$ avec $A, B \in \R[X]$~: alors
\[
P = (A + \iu B)(A - \iu B) = A^2 + B^2 .
\]
\end{solution}

\begin{solution}{exo:b1:poly:8}
Lagrange (\cref{thm:b1:poly:lagrange}) avec les nœuds $0, 1, 2$~:
\[
P = 1\cdot\frac{(X-1)(X-2)}{(0-1)(0-2)} + 3\cdot\frac{X(X-2)}{1\cdot(1-2)}
+ 2\cdot\frac{X(X-1)}{2\cdot 1}
= \frac{(X-1)(X-2)}{2} - 3X(X-2) + X(X-1).
\]
En développant~: $\frac{X^2 - 3X + 2}{2} - 3X^2 + 6X + X^2 - X =
-\frac{3}{2}X^2 + \frac{7}{2}X + 1$.

Le système~: $P = aX^2 + bX + c$ avec $c = 1$~; $a + b + 1 = 3$~; $4a +
2b + 1 = 2$. En retranchant deux fois la deuxième équation à la
troisième~: $2a - 1 = -4$, donc $a = -\frac32$, $b = \frac72$. Le même
\omterm{def:b1:poly:def}{polynôme}~: $P = -\frac32 X^2 + \frac72 X + 1$. (Vérification~: $P(2) =
-6 + 7 + 1 = 2$.)
\end{solution}

\begin{solution}{exo:b1:poly:9}
$P = X^{2n+1} - 1$~: $P' = (2n+1)X^{2n} \geq 0$, donc la fonction
polynomiale est croissante (strictement sauf en $0$), de limites
$\mp\infty$~: elle s'annule exactement une fois sur $\R$ (en $x = 1$).

Posons $E_n = \sum_{k=0}^{n} \frac{X^k}{k!}$. Alors $E_n' = E_{n-1} =
E_n - \frac{X^n}{n!}$. Une racine multiple $a$ vérifierait $E_n(a) =
E_n'(a) = 0$ (\cref{prop:b1:poly:multiplicity}), d'où
$\frac{a^n}{n!} = E_n(a) - E_n'(a) = 0$, donc $a = 0$~; or $E_n(0) = 1
\neq 0$. Aucune racine multiple.
\end{solution}

\begin{solution}{exo:b1:poly:10}
\begin{enumerate}
  \item Récurrence (les deux cas initiaux sont vrais). En utilisant
        $\cos(n+1)\theta + \cos(n-1)\theta = 2\cos\theta\cos
        n\theta$~:
        \[
        T_{n+1}(\cos\theta) = 2\cos\theta \cos n\theta -
        \cos(n-1)\theta = \cos(n+1)\theta .
        \]
  \item $T_n(\cos\theta) = 0$ si et seulement si $n\theta \equiv
        \frac\pi2 \pmod \pi$~: les nombres
        \[
        x_k = \cos\Bigl(\frac{(2k+1)\pi}{2n}\Bigr),
        \qquad k = 0, 1, \dots, n-1,
        \]
        sont $n$ points distincts de $\intoo{-1}{1}$ (les angles sont
        dans $\intoo{0}{\pi}$ où $\cos$ est \omterm{def:b1:logic:inj}{injective}), tous racines
        de $T_n$~; comme $\deg T_n = n$ (par la récurrence, avec
        coefficient dominant $2^{n-1}$ pour $n \geq 1$), ce sont
        \emph{toutes} les racines, chacune simple.
  \item Pour $x = \cos\theta \in \intcc{-1}{1}$~: $\abs{T_n(x)} =
        \abs{\cos n\theta} \leq 1$, avec égalité si et seulement si
        $n\theta \equiv 0 \pmod\pi$, c'est-à-dire aux $n+1$ points
        $y_k = \cos\frac{k\pi}{n}$, $k = 0, \dots, n$, où $T_n(y_k) =
        (-1)^k$. (C'est cette équioscillation qui fait de
        $2^{1-n}T_n$ le \omterm{def:b1:poly:def}{polynôme unitaire} de degré $n$ de plus petite
        norme uniforme sur $\intcc{-1}{1}$ --- démontré dans le
        devoir maison de ce chapitre.)
\end{enumerate}
\end{solution}

\begin{solution}{exo:b1:poly:11}
Écrivons $P = c\prod_i (X - a_i)^{m_i}$. La règle du produit (étendue à
plusieurs facteurs) donne
\[
P' = c\sum_{i} m_i (X - a_i)^{m_i - 1} \prod_{k \neq i} (X -
a_k)^{m_k},
\]
et en divisant par $P$~: $\frac{P'}{P} = \sum_i \frac{m_i}{X - a_i}$
(comme fractions rationnelles, c'est-à-dire hors des racines).

Soit $w$ une racine de $P'$. Si $w$ est l'un des $a_i$, il appartient
trivialement à l'enveloppe convexe. Sinon, en évaluant en $w$~:
\[
0 = \sum_i \frac{m_i}{w - a_i}
= \sum_i m_i\, \frac{\conj w - \conj a_i}{\abs{w - a_i}^2} .
\]
En conjuguant~: $\sum_i \lambda_i (w - a_i) = 0$ où $\lambda_i =
\frac{m_i}{\abs{w - a_i}^2} > 0$. D'où
\[
w = \frac{\sum_i \lambda_i a_i}{\sum_i \lambda_i} :
\]
une combinaison convexe (poids positifs de somme $1$ après
normalisation) des racines $a_i$. Ainsi toute racine de $P'$ appartient
à l'enveloppe convexe des racines de $P$.
\end{solution}

\begin{solution}{exo:b1:poly:12}
Sommons les évaluations de $(1 + X)^n$ aux trois racines cubiques de
l'unité~:
\[
2^n + (1 + j)^n + (1 + j^2)^n
= \sum_{k=0}^n \binom nk\,\bigl(1 + j^k + j^{2k}\bigr)
= 3\sum_{k\,:\,3\mid k}\binom nk ,
\]
puisque $1 + j^k + j^{2k}$ est une somme géométrique valant $3$ quand
$3 \mid k$ et $\frac{j^{3k} - 1}{j^k - 1} = 0$ sinon. Or $1 + j =
\frac12 + \iu\frac{\sqrt3}2 = \eu^{\iu\pi/3}$ et $1 + j^2 = \conj{1 +
j} = \eu^{-\iu\pi/3}$, donc $(1+j)^n + (1+j^2)^n = 2\cos\frac{n\pi}3$
et
\[
\sum_{k\geq0}\binom n{3k} = \frac{2^n + 2\cos\frac{n\pi}3}{3} .
\]
Vérifications~: $n = 3$~: $\frac{8 + 2\cos\pi}3 = 2 = \binom30 +
\binom33$~; $n = 6$~: $\frac{64 + 2}3 = 22 = 1 + 20 + 1$.
\end{solution}

\begin{solution}{pb:b1:poly:1}
\textbf{1.} $T_2 = 2X^2 - 1$~; $T_3 = 2X(2X^2 - 1) - X = 4X^3 -
3X$~; $T_4 = 2X\,T_3 - T_2 = 8X^4 - 8X^2 + 1$~; $T_5 = 2X\,T_4 - T_3
= 16X^5 - 20X^3 + 5X$. L'identité $T_3(\cos\theta) = \cos3\theta$
est exactement $\cos3\theta = 4\cos^3\theta - 3\cos\theta$ de
l'\cref{ex:b1:complex:cos3}.

\textbf{2.} Vrai pour $n = 1, 2$. Si $T_{n-1}$ et $T_n$ sont de
degrés $n-1$ et $n$, de coefficients dominants $2^{n-2}$ et
$2^{n-1}$, alors $2X\,T_n$ est de degré $n+1$ et de coefficient
dominant $2^n$, tandis que $T_{n-1}$ est de degré plus petit~:
$T_{n+1}$ est de degré $n + 1$ et de coefficient dominant $2^n$.
Parité~: si $T_{n-1}$ a la parité de $n - 1$ et $T_n$ celle de $n$,
alors $2X\,T_n$ et $T_{n-1}$ ont tous deux la parité de $n + 1$, donc
$T_{n+1}$ également.

\textbf{3.} Si $P(\cos\theta) = \cos n\theta$ pour tout $\theta$,
alors $P$ et $T_n$ coïncident en tout point de $\intcc{-1}1$ --- un
\omterm{def:b1:logic:sets}{ensemble} infini --- donc $P - T_n$ a une infinité de racines et est
le \omterm{def:b1:poly:def}{polynôme} nul (\cref{cor:b1:poly:nroots}).

\textbf{4.} Pour $x = \cos\theta$~: $T_m(T_n(\cos\theta)) =
T_m(\cos n\theta) = \cos(mn\theta) = T_{mn}(\cos\theta)$, et
$2T_mT_n(\cos\theta) = 2\cos m\theta\cos n\theta = \cos(m+n)\theta
+ \cos\abs{m - n}\theta$. Les deux identités valent sur
$\intcc{-1}1$, donc comme identités polynomiales par l'argument de
la question 3.

\textbf{5.} Les $x_k$ sont $n$ racines simples distinctes et le
coefficient dominant vaut $2^{n-1}$~:
\[
T_n = 2^{n-1}\prod_{k=0}^{n-1}
\Bigl(X - \cos\frac{(2k+1)\pi}{2n}\Bigr) .
\]
Entrelacement~: les angles $0 < \frac{\pi}{2n} < \frac\pi n <
\frac{3\pi}{2n} < \frac{2\pi}n < \dots < \pi$ alternent entre les
angles $\frac{k\pi}n$ des $y$ et les angles $\frac{(2k+1)\pi}{2n}$
des $x$~; comme $\cos$ est strictement décroissante sur
$\intcc0\pi$, les valeurs s'entrelacent dans l'ordre inverse~: $y_n
< x_{n-1} < y_{n-1} < \dots < x_0 < y_0$. Entre deux extrema
consécutifs se trouve exactement une racine, comme le suggère le
graphe de $\cos n\theta$.

\textbf{6.} Récurrence avec $2\cosh a\cosh b = \cosh(a + b) +
\cosh(a - b)$ (\cref{prop:b1:functions:hyprules})~: $T_{n+1}(\cosh
t) = 2\cosh t\cosh nt - \cosh(n-1)t = \cosh(n+1)t$. Pour $x \geq
1$, écrivons $x = \cosh t$ avec $t \geq 0$~; alors $\eu^t = x +
\sqrt{x^2 - 1}$ et $\eu^{-t} = x - \sqrt{x^2 - 1}$, donc
\[
T_n(x) = \cosh(nt)
= \frac{(x + \sqrt{x^2-1})^n + (x - \sqrt{x^2-1})^n}2 .
\]
Pour $x > 1$, le premier terme dépasse strictement $\frac12(1)^n$ et
croît géométriquement~: $T_n(x) > 1$.

\textbf{7.} De Moivre~: $\cos n\theta = \Re\bigl((\cos\theta +
\iu\sin\theta)^n\bigr) = \sum_{2j \leq n}\binom n{2j}
\cos^{n-2j}\theta\,(\iu\sin\theta)^{2j}$, et $(\iu\sin\theta)^{2j}
= (-\sin^2\theta)^j = (\cos^2\theta - 1)^j$. En substituant $x =
\cos\theta$ et en invoquant la question 3~:
\[
T_n(x) = \sum_{0\leq 2j\leq n}\binom n{2j}x^{n-2j}(x^2 - 1)^j .
\]
Pour $n = 3$~: $\binom30 x^3 + \binom32 x(x^2 - 1) = x^3 + 3x^3 -
3x = 4x^3 - 3x$, comme à la question 1.

\textbf{8.} $T_n(1) = \cos(n\cdot0) = 1$~; $T_n(-1) = \cos(n\pi) =
(-1)^n$~; $T_n(0) = \cos\frac{n\pi}2$, qui vaut $0$ pour $n$ impair
et $(-1)^{n/2}$ pour $n$ pair.

\textbf{9.} Récurrence pour $U_n(\cos\theta) =
\frac{\sin(n+1)\theta}{\sin\theta}$~: vrai pour $U_0 = 1$ et $U_1 =
2X$ ($\sin2\theta = 2\sin\theta\cos\theta$)~; l'hérédité est
l'identité de transformation $\sin(n+2)\theta = 2\cos\theta\,
\sin(n+1)\theta - \sin n\theta$. Dérivons maintenant
$T_n(\cos\theta) = \cos n\theta$ par rapport à $\theta$~:
$-\sin\theta\,T_n'(\cos\theta) = -n\sin n\theta$, donc pour $\theta
\notin \pi\Z$~:
\[
T_n'(\cos\theta) = n\,\frac{\sin n\theta}{\sin\theta}
= n\,U_{n-1}(\cos\theta) ,
\]
et les \omterm{def:b1:poly:def}{polynômes} $T_n'$ et $nU_{n-1}$, qui coïncident sur
$\intoo{-1}1$, sont égaux.

\textbf{10.} $\abs{\sin(n+1)\theta} = \abs{\sin n\theta\cos\theta
+ \cos n\theta\sin\theta} \leq \abs{\sin n\theta} +
\abs{\sin\theta}$, et la récurrence donne $\abs{\sin n\theta} \leq
n\abs{\sin\theta}$. D'où $\abs{U_{n-1}} \leq n$ sur $\intoo{-1}1$
et $\abs{T_n'} = n\abs{U_{n-1}} \leq n^2$ sur cet intervalle~; en
$\pm1$, la majoration s'étend par passage à la limite (ou
directement~: $U_{n-1}(1) = n$ par la récurrence, $U_n(1) = n + 1$
par récurrence, et la parité donne $U_{n-1}(-1) = (-1)^{n-1}n$).
Ainsi $T_n'(1) = n^2$ et $T_n'(-1) = (-1)^{n-1}n^2$~: la borne $n^2$
est atteinte aux extrémités.

\textbf{11.} Dérivons $\sin\theta\,T_n'(\cos\theta) = n\sin
n\theta$ (question 9) par rapport à $\theta$~:
\[
\cos\theta\,T_n'(\cos\theta) - \sin^2\theta\,T_n''(\cos\theta)
= n^2\cos n\theta = n^2\,T_n(\cos\theta) .
\]
Avec $x = \cos\theta$ et $\sin^2\theta = 1 - x^2$~: $x\,T_n' - (1
- x^2)T_n'' = n^2T_n$ sur $\intcc{-1}1$, donc partout~:
$(1 - x^2)y'' - xy' + n^2y = 0$ pour $y = T_n$. Vérification pour
$T_2 = 2x^2 - 1$~: $(1 - x^2)(4) - x(4x) + 4(2x^2 - 1) = 4 - 4x^2 -
4x^2 + 8x^2 - 4 = 0$.

\textbf{12.} $\widetilde T_n$ est \omterm{def:b1:poly:def}{unitaire} (question 2) et
$\abs{\widetilde T_n} = 2^{1-n}\abs{T_n} \leq 2^{1-n}$ sur
$\intcc{-1}1$, avec $\widetilde T_n(y_k) = (-1)^k2^{1-n}$ aux $n +
1$ points $y_k$ (\cref{exo:b1:poly:10})~: la norme vaut exactement
$2^{1-n}$, atteinte avec des signes alternés.

\textbf{13.} $\widetilde T_n$ et $P$ sont tous deux \omterm{def:b1:poly:def}{unitaires} de
degré $n$, donc les termes dominants se simplifient~: $\deg D \leq n
- 1$. En $y_k$~: $D(y_k) = (-1)^k2^{1-n} - P(y_k)$, et $\abs{P(y_k)}
\leq \norm P_\infty < 2^{1-n}$ impose au signe de $D(y_k)$ d'être
celui de $(-1)^k2^{1-n}$, strictement.

\textbf{14.} $D$ change de signe entre $y_{k+1}$ et $y_k$ pour
chaque $k = 0, \dots, n-1$~: par la propriété des valeurs
intermédiaires, $D$ a une racine dans chacun de ces $n$ intervalles
ouverts deux à deux disjoints --- $n$ racines distinctes pour un
\omterm{def:b1:poly:def}{polynôme} non nul de degré $\leq n - 1$, impossible. Et $D = 0$ est
également impossible (les normes diffèrent). Contradiction~: aucun
\omterm{def:b1:poly:def}{polynôme unitaire} $P$ de degré $n$ ne vérifie $\norm P_\infty <
2^{1-n}$, ce qui est le théorème de Tchebychev.

\textbf{15.} Maintenant $\abs{P(y_k)} \leq 2^{1-n}$ seulement, donc
$(-1)^k D(y_k) = 2^{1-n} - (-1)^kP(y_k) \geq 2^{1-n} - \abs{P(y_k)}
\geq 0$. Supposons $D(y_k) = 0$ en un $y_k$ intérieur ($0 < k <
n$)~: alors $P(y_k) = (-1)^k2^{1-n}$, donc $\abs P$ atteint sa borne
supérieure $2^{1-n}$ au point intérieur $y_k$, d'où $P'(y_k) = 0$
(extremum intérieur)~; et $T_n'(y_k) = nU_{n-1}(y_k) = 0$ puisque
$\sin(n\cdot\frac{k\pi}n) = 0$ --- donc $\widetilde T_n'(y_k) = 0$
aussi, et $D'(y_k) = 0$~: $y_k$ est racine de $D$ de \omterm{def:b1:poly:derivative}{multiplicité} au
moins $2$.

\textbf{16.} Comptons les racines de $D$ avec \omterm{def:b1:poly:derivative}{multiplicité}. Soit $z$
le nombre de points intérieurs $y_k$ tels que $D(y_k) = 0$ (chacun
racine double, par la question 15) et $e \in \{0, 1, 2\}$ le nombre
d'extrémités ($y_0$ ou $y_n$) où $D = 0$ (chacune racine au moins
simple). Un intervalle $(y_{k+1}, y_k)$ dont les deux extrémités
vérifient $D \neq 0$ porte des signes strictement alternés, donc une
racine intérieure. Chaque point intérieur annulé gâte au plus les
deux intervalles adjacents, chaque extrémité annulée au plus un
intervalle~: au moins $n - 2z - e$ intervalles fournissent encore une
racine chacun, toutes distinctes des racines en les $y$. Total~: au
moins $(n - 2z - e) + 2z + e = n$ racines comptées avec \omterm{def:b1:poly:derivative}{multiplicité},
pour un \omterm{def:b1:poly:def}{polynôme} de degré $\leq n - 1$~: donc $D = 0$ et $P =
\widetilde T_n$. Le minimiseur est unique.

\textbf{17.} L'\omterm{def:b1:logic:map}{application} affine $t \mapsto x = \frac{a+b}2 +
\frac{b-a}2\,t$ est une \omterm{def:b1:logic:inj}{bijection} de $\intcc{-1}1$ sur $\intcc ab$.
Si $P$ est \omterm{def:b1:poly:def}{unitaire} de degré $n$, alors $Q(t) = P(x(t))$ est un
\omterm{def:b1:poly:def}{polynôme} en $t$ de coefficient dominant
$\bigl(\frac{b-a}2\bigr)^n$, et $\sup_{\intcc ab}\abs P =
\sup_{\intcc{-1}1}\abs Q$. Le \omterm{def:b1:poly:def}{polynôme unitaire}
$Q/\bigl(\frac{b-a}2\bigr)^n$ a une norme uniforme $\geq 2^{1-n}$
(questions 13--14), donc
\[
\sup_{\intcc ab}\abs P \geq \Bigl(\frac{b-a}2\Bigr)^n 2^{1-n}
= 2\Bigl(\frac{b-a}4\Bigr)^n ,
\]
avec égalité exactement pour $P(x) = \bigl(\frac{b-a}2\bigr)^n
\widetilde T_n\bigl(t(x)\bigr)$ (question 16).

\textbf{18.} $\widetilde T_3 = \frac{T_3}4 = X^3 - \frac34X$~;
$\widetilde T_3{}' = 3X^2 - \frac34$ s'annule en $\pm\frac12$.
Valeurs~: $\widetilde T_3(-1) = -\frac14$, $\widetilde
T_3(-\tfrac12) = \frac14$, $\widetilde T_3(\tfrac12) = -\frac14$,
$\widetilde T_3(1) = \frac14$~: quatre extrema alternés de valeur
absolue $\frac14$ --- donc $\norm{\widetilde T_3}_\infty =
\frac14$, et par le théorème de Tchebychev aucune cubique \omterm{def:b1:poly:def}{unitaire}
n'a une norme uniforme plus petite sur $\intcc{-1}1$.

\textbf{19.} $\omega$ est \omterm{def:b1:poly:def}{unitaire} de degré $n + 1$, donc
$\norm\omega_\infty \geq 2^{-n}$ par le théorème de Tchebychev (en
degré $n+1$), avec égalité si et seulement si $\omega = \widetilde
T_{n+1} = 2^{-n}T_{n+1}$ (question 16), c'est-à-dire si et seulement
si les nœuds sont les $n + 1$ racines de $T_{n+1}$. Avec les nœuds
de Tchebychev, le facteur d'erreur $\norm\omega_\infty$ vaut $2^{-n}$
--- le plus petit possible.

\textbf{20.} $5 \times 36^\circ = 180^\circ$, donc $T_5(c) =
\cos180^\circ = -1$~: $16c^5 - 20c^3 + 5c + 1 = 0$. En testant $x =
-1$~: $-16 + 20 - 5 + 1 = 0$, et le développement confirme
\[
16x^5 - 20x^3 + 5x + 1 = (x + 1)\bigl(4x^2 - 2x - 1\bigr)^2 .
\]
Comme $c = \cos36^\circ \neq -1$, $c$ est racine de $4x^2 - 2x -
1$, dont les racines sont $\frac{1 \pm \sqrt5}4$~; comme $c > 0$,
\[
\cos36^\circ = \frac{1 + \sqrt5}4 .
\]
Cohérence~: $\cos72^\circ = T_2(c) = 2c^2 - 1 = 2\cdot\frac{3 +
\sqrt5}8 - 1 = \frac{\sqrt5 - 1}4$, la valeur trouvée dans
l'\cref{exo:b1:complex:8}.

\textbf{21.} $\sqrt{1.1^2 - 1} = \sqrt{0.21} \approx 0.458$, donc
$x + \sqrt{x^2-1} \approx 1.558$ et $(1.558)^{10} \approx 84.5$,
tandis que $(1.1 - 0.458)^{10} \approx 0.01$~: $T_{10}(1.1) \approx
\frac{84.5 + 0.01}2 \approx 42$. Un \omterm{def:b1:poly:def}{polynôme} confiné dans
$\intcc{-1}1$ sur l'intervalle a déjà dépassé $40$ un dixième
au-delà du bord~: être borné sur un segment ne dit rien à un pouce
de là.

\textbf{22.} Dans la formule de la question 7 pour $T_p$, le terme
$j = 0$ vaut $X^p$~; tout autre terme porte $\binom p{2j}$ avec $0 <
2j < p$ (noter que $2j \neq p$ puisque $p$ est impair), qui est
divisible par $p$ d'après la première étape de la démonstration du
\cref{thm:b1:arith:fermat}. Donc tout coefficient de $T_p - X^p$ est
un multiple de $p$. Vérifications~: $T_3 - X^3 = 3X^3 - 3X = 3(X^3 -
X)$~; $T_5 - X^5 = 15X^5 - 20X^3 + 5X = 5(3X^5 - 4X^3 + X)$.

\textbf{23.} Par la question 17 avec $\intcc ab = \intcc01$ et $n =
2$~: écart minimal $2\bigl(\frac14\bigr)^2 = \frac18$, atteint par
$\bigl(\frac12\bigr)^2\widetilde T_2(2x - 1) =
\frac14\bigl((2x-1)^2 - \frac12\bigr) = x^2 - x + \frac18$. Le
trinôme \omterm{def:b1:poly:def}{unitaire} le plus proche de zéro sur $\intcc01$ est $x^2 - x
+ \frac18$, de norme uniforme $\frac18$.

\textbf{24.} (i) La rigidité --- un \omterm{def:b1:poly:def}{polynôme} ayant plus de racines
que son degré est nul --- a alimenté le principe d'unicité
(question 3), le transfert des identités trigonométriques en
identités polynomiales (questions 4, 7, 9, 11), et les deux
décomptes de racines de la démonstration d'extrémalité (questions
14, 16). (ii) La trigonométrie du \cref{ch:b1:complex} (de Moivre,
transformation de produits en sommes) et les \omterm{def:b1:functions:hyperbolic}{fonctions hyperboliques}
du \cref{ch:b1:functions} ont fourni toutes les identités qui
soutiennent la famille~; la substitution $x = \cos\theta$ en est le
pont. (iii) La \omterm{def:b1:arith:divides}{divisibilité} $p \mid \binom p{2j}$ du
\cref{ch:b1:arith} a transformé la formule des coefficients en la
\omterm{def:b1:arith:congruence}{congruence} de la question 22.

\textbf{25.} Le théorème de Tchebychev convertit une optimisation
sur une famille de dimension infinie (tous les \omterm{def:b1:poly:def}{polynômes unitaires})
en une combinatoire finie~: un concurrent meilleur que $\widetilde
T_n$ en différerait par un \omterm{def:b1:poly:def}{polynôme} de bas degré contraint de
changer $n$ fois de signe --- une racine de plus que son degré ne le
permet. Le motif d'équioscillation n'est donc pas une curiosité,
mais le certificat même de l'optimalité, et le cas d'égalité affine
le décompte des racines en tenant compte des \omterm{def:b1:poly:derivative}{multiplicités}. La
substitution $x = \cos\theta$ mérite le dernier mot~: elle
transporte le monde rigide et discret des \omterm{def:b1:poly:def}{polynômes} dans le monde
périodique de la trigonométrie, où racines et extrema de $T_n$ ne
sont que la grille régulière de $\cos n\theta$. Les deux résultats
d'analyse empruntés --- la propriété des valeurs intermédiaires
(question 14~; démontrée au \cref{ch:b1:continuity}) et la nullité
de la dérivée en un extremum intérieur (question 15~; démontrée au
\cref{ch:b1:derivative}) --- sont exactement les outils que ces
chapitres ultérieurs rendront, bouclant la boucle.
\end{solution}
